\section{Related Work}
\label{sec:related}

\subsection{Object Detection for Small Targets}
The one-stage and two-stage families are both represented in our benchmark: Faster
R-CNN~\cite{ren2015faster} as the two-stage baseline, ATSS~\cite{zhang2020atss} and
TOOD~\cite{feng2021tood} as anchor-assignment refinements, RTMDet~\cite{lyu2022rtmdet} as a
modern convolutional one-stage design, and the YOLO
line~\cite{redmon2016yolo,wang2024yolov9,wang2024yolov10} as the practical default.
Query-based detectors --- DETR~\cite{carion2020detr}, its accelerated variants, and
DINO~\cite{zhang2023dino} --- are known to converge slowly and to struggle on very small
objects, and RT-DETR~\cite{zhao2024rtdetr} targets that gap. Our animals are a median
$29\times24$ pixels, which is squarely the regime where these differences matter, and
\cref{sec:detection} reports how each family transfers.

\subsection{Multi-Object Tracking}
Tracking-by-detection remains dominant: SORT~\cite{bewley2016sort} and
DeepSORT~\cite{wojke2017deepsort} established the Kalman-plus-association template,
ByteTrack~\cite{zhang2022bytetrack} showed that associating \emph{low}-confidence detections
recovers substantial recall, and BoT-SORT~\cite{aharon2022botsort} added camera-motion
compensation and an appearance term. Our camera is mounted on a vehicle moving at
25--90\,km/h, so global motion compensation is not optional; we use BoT-SORT with sparse
optical-flow compensation throughout and ablate its appearance term in
\cref{sec:inversion}. Critically, MOT benchmarks score association quality (MOTA, IDF1,
HOTA), not the number of \emph{distinct objects that ever existed} --- which is the quantity
a survey needs, and the quantity our decomposition measures.

\subsection{Animal Detection and Counting}
Camera-trap classification at scale~\cite{norouzzadeh2018automatically} and the domain-shift
benchmarks that followed~\cite{beery2018recognition} established animal imagery as its own
distribution problem, while aerial and UAV surveys have received the most counting-specific
attention~\cite{kellenberger2018detecting,corcoran2019automated}, usually with per-image
counts on sparse stills rather than per-individual identity across continuous video.
Thermal wildlife datasets with \emph{per-individual track} annotation over long sequences
remain scarce; that scarcity is what our corpus addresses.

\subsection{Counting as Its Own Task}
Crowd counting~\cite{zhang2016single} shows what a field looks like when it takes counting
seriously as a distinct objective --- density regression rather than detection-then-count.
That approach suits dense scenes with hundreds of overlapping instances; ours is the
opposite regime, sparse instances that must each be resolved as an individual because
survey statistics are computed per animal. Our contribution is to show that in this sparse
regime, the detection-then-count framing is not wrong so much as \emph{mis-attributed}: the
detector is already near-saturated, and the accuracy that a practitioner would try to buy
with a better backbone is instead being lost downstream.

\subsection{Animal Re-Identification}
Individual identification from appearance is established for species with high-contrast,
high-resolution natural markings. \Cref{sec:reid} tests whether it transfers here: at a
median 27-pixel animal, appearance carries less identity signal than a two-number size
descriptor, and an ImageNet backbone scores \emph{below} plain geometry.
